%%%% For glossarie


%%%%%%%%%%%%%%%%%%%%
%%%%%% Biochemestry part
%%%%%%%%%%%%%%%%%%

\newglossaryentry{archa}{
	name={archaea},
	description={Archaea constitute one of the three major domains of life (bacteria, archaea, eukaryotes), and, like bacteria, they belong to the prokaryotes \citep[p. \,3]{fritscheMikrobiologie2024}},
	type=biochem
}

\newglossaryentry{suicide}{name=suicide plasmid,description={A plasmid that has no \gls{ori} and therefore can not be replicated in the cell except it gets integrated in the \acrshort{gDna}}, 
type=biochem}

\newglossaryentry{auxo}{name=auxotrophic,description= {Auxotrophic organism can not produce relevant bio-molecules for growth and surviving. Therefor, they must obtain these from their environment \parencite[p.\,65]{fritscheMikrobiologie2024}}, 
type=biochem}

\newglossaryentry{complex}{
	name={complex media},
	description={Complex media contain natrual, organic moleculas, that are not precisely defined chemically(e.g. yeast extract or pepton) So a lot of diffrent microbes can grow on them \citep[p. \,359]{fritscheMikrobiologie2024}},
	type=biochem
}

\newglossaryentry{pha_con_mic}{
	name={phase-contrast microscope},
	description={Phase-contrast microscopy is based on the fact, that objects or medium not only changes the amplitude of light but also the phase of it. The phase-contrast microscope makes this phenomena visual for the human eye \citep{Phasecontrast_uni_leipzig_2009}},
	type=biochem
}

\newglossaryentry{plateing}{
	name={Plating},
	description= {Cultivate microorganism on a solid medium (contains a gelling agent like agar). THerby is the plating the process how the organism get distributed on the plate / solid media. Different technics are possible like spreading or pouring \citep[p. \,364]{fritscheMikrobiologie2024}},
	type=biochem}

\newglossaryentry{pour}{
	name={pour plating},
	description= {A plating technique in which the solution containing microorganisms is mixed with the agar-containing culture medium while it is still liquid (at about 45 °C). This mixture is then poured into a Petri dish, where it solidifies\citep[p. \,364]{fritscheMikrobiologie2024}},
	type=biochem}
	
\newglossaryentry{spread}{
	name={spread plating},
	description= {A plating technique in which the solution containing microorganisms is put on the solid agar-containing culture medium. The solution is then spread across the agar plate using a Dragalski spatula \citep[p. \,364]{fritscheMikrobiologie2024}},
	type=biochem}

\newglossaryentry{cterm}{name={C-terminal end},
	description={The end of the amino acid sequence of a protein, where the carboxy group of the terminal amino acid has no covalent connection to another amino acid},
	type=biochem}

\newglossaryentry{nterm}{name={N-terminal end},
	description={The end of the amino acid sequence of a protein, where the amino group of the terminal amino acid has no covalent connection to another amino acid},
	type=biochem}

\newglossaryentry{promo}{
name=promotor,description= {Region on DNA, which allows binding of relevant enzyms for the transcription of genes. So it is a essential region for gene expression}, 
type=biochem}

\newglossaryentry{trafoglo}{name=transformation,description= {The process of getting a extern DNA into the target organism and its expression. If the target organism is a animal eucaryotic cell it is called transfection}, 
type=biochem}

\newglossaryentry{ori}{name = ORI, description =  {Origin of replication, is the region of the DNA, where the machinery responsible for replicating the entire genomic DNA first binds and begins replication}, 
type = biochem}

\newglossaryentry{plasm}{
name=plasmid,
description={ A small ring of double-stranded DNA. Occurs naturally in \gls{archa} and bacteria, but is also used for scientific purposes. Here often as a vector for a transformation},
type = biochem}

\newglossaryentry{alphafold}{name=AlphaFold,description={Is a KI based software for prediction of 3-D structure of proteins. For the prediction it needs an aminoacid sequence}, 
type=biochem}

\newglossaryentry{rstrsite}{name=restricitionsite,description={Restrictionsites are a specific sequence of approx. 4 to 10 bp that are recognized by a specific restriction-enzyme. This cuts the double-stranded DNA at this site in a specific behavior}, 
type=biochem }

\newglossaryentry{betaFaltblatt}{name=$\beta$-sheet, 
description={A common secondary structure of proteins. It consists of $\beta$-strands arranged in parallel or antiparallel orientations. $\beta$-strands are illustrated as arrows. Another typical secondary structure is the $\alpha$-helix}, 
type=biochem}

\newglossaryentry{glove}{
	name={glovebox},
	description={A container that is gas-tight, so that a different atmosphere can be created},
	type=biochem}

\newglossaryentry{col}{
	name= {colony},
	description={In a microbiological context, a colony is defined as a clearly distinct cluster of cells that can be traced back to a single parent cell through cell division}, 
	type=biochem}

\newglossaryentry{ribo}{
	name={ribosome},
	description={Ribosomes are a mix out \acrfull{rrna} and proteins. They play a central role in protein biosynthesis by catalyzing the translation of mRNA into an amino acid sequence},
	type=biochem}


\newglossaryentry{S-sequenz}{
	name={16S-Sequencing},
	description={16S sequencing targets the 16s genes, which code for the 16S \acrfull{rrna}. The 16S rRNA is a part of the procaryotic \gls{ribo}. These 16-s base sequenz is specific for the organism},
	type=biochem
}

\newglossaryentry{primer}{
	name=primer,
	description={Short nucleotide sequence that serves as a starting point for the polymerase to replicate the complementary DNA strand. For this purpose, the primer hybridizes to a specific sequence on the \acrfull{ssDNA} template \citep{fritscheMikrobiologie2024}},
	type=biochem
}

\newglossaryentry{invitro}{name=in-vitro,description= {Experiments are not performed in/ at the living organism, but in an artificial setting}, 
type=biochem}

\newglossaryentry{pcr}{
	name=PCR,
	description={The polymerase chain reaction is a \gls{invitro} method to amplify a specific sequence of DNA. For this purpose, specific fitting \gls{primer} has to be designed, that flank this DNA-sequence on both complementary DNA-strings.}, 
	type=biochem}

%%%%%%%%%%%%%%%%%%%%
%%%%%% Geography part
%%%%%%%%%%%%%%%%%%


\newglossaryentry{crtglo}{name=CRT,description={Common reporting table is a template for countries to report its emission of GHG fitting IPCC standards. The UNCFF has oblige the ANNEX-I countries to publish this every year since 1990}, 
type=geo}

\newglossaryentry{annex_one}{name=ANNEX-I,
description={ANNEX-I and Non-ANNEX-I countries is a classification of the \acrshort{unfccc}. ANNEX-I are the \acrshort{oecd} countries or with an economy in transition}, 
type=geo}

\newglossaryentry{gwpGls}{name=GWP,description={Global warming potential is a concept used to represent the relative impact of various \acrshort{ghg}s on a uniform scale over a specific time period. It takes into account their effectiveness in causing radiative forcing as well as their atmospheric lifetime. \ch{CO2} serves as a reference, so that the effect of 1 g of greenhouse gas is converted to x g of \ch{CO2} \acrshort{eq}. The relevant time period is indicated as the number of years following the GWP designation (e. g., GWP20 refers to a 20-year time horizon)}, 
type=geo}

\newglossaryentry{compound}{name=compound events,description={When two wether extremes occur at the same time a shortly after each other, so so that they reinforce each other, as example a heatwave and droughts \citep{leeAnnexGlossary2023}}, 
type=geo }

\newglossaryentry{tippingPoint}{name=tipping point, description={"Refers to a critical point at which small changes can suddenly trigger large, often irreversible  shifts in Earth’s environment. Once a tipping point is crossed, self-reinforcing (positive feedback) processes drive the system further away from its previous state, increasing the magnitude and extent of change. For example, melting ice exposes less reflective ocean water, which absorbs more sunlight and accelerates melting, creating a self-reinforcing cycle" \citep{planetaryboundariessciencepbsciencePlanetaryHealthCheck2025}}, 
type=geo}

\newglossaryentry{zoneOfIncreasingRisk}{name=zone of an increasing risk,description= {"Refers to the stage when human activities push Earth beyond the Planetary Boundaries, entering a “Zone of Increasing Risk.” In this zone, the further boundaries are exceeded, the greater the chance of causing serious damage, destabilizing key Earth system processes, and disrupting life-support functions" \citep{planetaryboundariessciencepbsciencePlanetaryHealthCheck2025}}, 
type=geo}

\newglossaryentry{HighRiskZone}{name= high-risk zone,description={"Refers to Earth entering the “High-Risk Zone,” where there is a strong possibility of severe,  irreversible damage to key planetary functions that support life. In this zone, immediate action becomes critical to prevent locking in permanent changes and moving even further away from the stable conditions of the Holocene epoch (a period of stability on Earth covering approximately the last 12,000 years)" \parencite{planetaryboundariessciencepbsciencePlanetaryHealthCheck2025}}, 
type=geo}

\newglossaryentry{planetaryBoundarysGLS}{name=planetary boundaries,description={Is a concept that uses the variability of values of the Holocene as a reference to evaluate thresholds for different processes to maintain a stable system on earth, where human mankind can definitely thrive. The nine main processes are "Climate change", "Biosphere integrity", "Land-system change", "Freshwater use", "Biogeochemical flows", "Ocean acidification", "Stratospheric ozone depletion", "Atmospheric aerosol loading" and "Novel entities" \parencite{planetaryboundariessciencepbsciencePlanetaryHealthCheck2025}}, 
type=geo}


\newglossaryentry{radForce}{name=radiative forcing,description={Is a concept that quantifies the change of the energy balance of the hole planet. The unit is W/m\textsuperscript{2}. It depends on many factors as e.g. aeorosol conten of the atmosphere, planetary albedo and the concentration of \acrshort{ghg} \parencite{Chapter Eight - Radiative forcing on climate change: assessing the effect of greenhouse gases on energy balance of Earth}}, 
type=geo}
%%%%%%%%%%%%%%%%%%%"%%%%%%%%%%%%%%%%%%%%%%%%%%%%%%%%%%%%%%%%%%%%%

%%%%%%%%%%%%%%%%%%%"%%%%%%%%%%%%%%%%%%%%%%%%%%%%%%%%%%%%%%%%%%%%%


\printglossary[type=biochem, title={BIOCHEMICAL TERMS}, nonumberlist]
\newpage
\printglossary[type=geo, title={GEOGRAPHICAL TERMS}, nonumberlist]

